\chapter{نتایج و تحلیل مقایسه‌ای}

\section{آمار گراف دانش و نمایه‌سازی برداری}

برای توصیف اندازه و پوشش منبع دانش، آمار گراف \lr{PostgreSQL} و مجموعه‌های برداری \lr{Milvus} از پایگاه مورد استفاده استخراج شد. گراف شامل ۳۲۸۰۸ موجودیت، ۶۲۰۵۸ واقعیت و ۱۰۲۸۲ قطعه متنی است. جدول زیر توزیع موجودیت‌ها را بر اساس نوع نشان می‌دهد.

\begin{table}[H]
\centering
\caption{توزیع موجودیت‌های گراف دانش بر اساس نوع}
\label{tab:graph-entity-types}
\begin{tabular}{lr}
\hline
نوع موجودیت & تعداد \\
\hline
نشانه & ۶۷۳۵ \\
مفهوم & ۶۵۳۱ \\
ملاک تشخیصی & ۴۹۴۲ \\
مشخصه & ۳۰۸۹ \\
مدت & ۲۵۹۹ \\
موجودیت حل‌نشده & ۲۱۹۲ \\
کد تشخیصی & ۲۰۷۵ \\
اختلال & ۱۶۴۲ \\
بیمار یا گروه هدف & ۱۱۱۴ \\
برچسب نمونه & ۹۲۹ \\
رفتار & ۷۰۷ \\
عنوان بخش & ۲۵۳ \\
\hline
جمع & ۳۲۸۰۸ \\
\hline
\end{tabular}
\end{table}

در جدول \lr{facts} تعداد ۶۲۰۵۸ رابطه ذخیره شده است و جدول \lr{mentions} ارتباط این واقعیت‌ها با قطعه‌های متن را نگهداری می‌کند. در نتیجه، هر ۱۰۲۸۲ قطعه متن دست‌کم از طریق یک اشاره به واقعیت‌های گراف متصل است.

\begin{table}[H]
\centering
\caption{آمار مجموعه‌های برداری \lr{Milvus}}
\label{tab:milvus-collections}
\begin{tabular}{lrr}
\hline
مجموعه برداری & تعداد بردار & بُعد بردار \\
\hline
موجودیت‌های نهایی \lr{canonical-seed-entities-v1} & ۱۳۰۸۲ & ۱۲۸ \\
واقعیت‌ها \lr{raw-triples-v1} & ۴۱۶۵۴ & ۲۵۶ \\
قطعه‌های منبع \lr{source-chunks-v1} & ۱۰۲۸۲ & ۱۰۲۴ \\
\hline
\end{tabular}
\end{table}

بردارهای موجودیت‌های نهایی، رابطه‌های نهایی، واقعیت‌ها و قطعه‌های منبع با مدل \lr{jina-embeddings-v4} تولید شده‌اند. شناسه موجودیت‌ها و قطعه‌ها میان \lr{PostgreSQL} و \lr{Milvus} مشترک است؛ بنابراین، نتیجه جست‌وجوی برداری مستقیماً به رکورد گراف یا متن منبع متناظر متصل می‌شود.

\section{مقایسه نتایج پیش‌بینی در مجموعه‌داده‌ها}

نتایج پیش‌بینی پنج روش در دو مجموعه‌داده در جدول‌های زیر گزارش شده است. جدول \ref{tab:results-dr} نتایج کلان مجموعه‌داده \lr{DR} و جدول \ref{tab:results-swmh} نتایج کلان مجموعه‌داده \lr{SWMH} را نشان می‌دهند. برای هر روش، دقت کلی و میانگین معیارهای صحت، بازخوانی و \lr{F1-score} به‌صورت کلان ارائه شده‌اند. میانگین کلان با وزن برابر برای همه برچسب‌ها محاسبه شده است؛ بنابراین، اثر برچسب‌های کم‌نمونه نیز در مقایسه حفظ می‌شود.

\begin{table}[H]
\centering
\small
\caption{مقایسه نتایج پیش‌بینی روش‌ها در مجموعه‌داده \lr{DR}}
\label{tab:results-dr}
\begin{tabular}{lrrrrr}
\hline
روش & تعداد & دقت کلی & میانگین صحت & میانگین بازخوانی & میانگین \lr{F1-score} \\
\hline
\lr{Vanilla RAG} & ۷۰ & \lr{68.57\%} & \lr{70.81\%} & \lr{77.03\%} & \lr{67.26\%} \\
\lr{HippoRAG 2} & ۷۰ & \lr{70.00\%} & \lr{71.46\%} & \lr{77.99\%} & \lr{68.56\%} \\
\lr{CoT} & ۷۰ & \lr{80.00\%} & \lr{77.08\%} & \lr{84.72\%} & \lr{77.81\%} \\
\lr{ToG-1} & ۷۰ & \lr{75.71\%} & \lr{74.41\%} & \lr{81.84\%} & \lr{73.78\%} \\
\lr{ToG-2} & ۷۰ & \lr{75.71\%} & \lr{74.41\%} & \lr{81.84\%} & \lr{73.78\%} \\
\hline
\end{tabular}
\end{table}

\begin{table}[H]
\centering
\small
\caption{مقایسه نتایج پیش‌بینی روش‌ها در مجموعه‌داده \lr{SWMH}}
\label{tab:results-swmh}
\begin{tabular}{lrrrrr}
\hline
روش & تعداد & دقت کلی & میانگین صحت & میانگین بازخوانی & میانگین \lr{F1-score} \\
\hline
\lr{Vanilla RAG} & ۶۵ & \lr{66.15\%} & \lr{78.49\%} & \lr{65.81\%} & \lr{64.80\%} \\
\lr{HippoRAG 2} & ۶۵ & \lr{64.62\%} & \lr{74.22\%} & \lr{64.12\%} & \lr{65.43\%} \\
\lr{CoT} & ۶۵ & \lr{70.77\%} & \lr{75.06\%} & \lr{70.53\%} & \lr{71.03\%} \\
\lr{ToG-1} & ۶۵ & \lr{70.77\%} & \lr{74.31\%} & \lr{70.64\%} & \lr{71.14\%} \\
\lr{ToG-2} & ۶۵ & \lr{82.54\%} & \lr{82.66\%} & \lr{78.59\%} & \lr{80.31\%} \\
\hline
\end{tabular}
\end{table}

در مجموع دو مجموعه‌داده، روش \lr{ToG-2} بالاترین عملکرد کلی را نشان می‌دهد و پس از آن \lr{CoT} و \lr{ToG-1} قرار می‌گیرند. با این حال، عملکرد روش‌ها میان دو مجموعه‌داده یکسان نیست؛ \lr{CoT} در \lr{DR} بهترین نتیجه را دارد، درحالی‌که \lr{ToG-2} در \lr{SWMH} برتری محسوسی نشان می‌دهد. این تفاوت در بخش تحلیل توضیح‌پذیری و هزینه اجرا نیز بررسی خواهد شد.

برای مشاهده رفتار روش‌ها در هر برچسب، معیارهای تفکیک‌شده در جدول‌های \ref{tab:per-label-dr} و \ref{tab:per-label-swmh} آمده‌اند؛ جدول نخست مربوط به \lr{DR} و جدول دوم مربوط به \lr{SWMH} است. ستون پشتیبانی، تعداد نمونه‌های واقعی هر برچسب را نشان می‌دهد.

\begin{table}[ht]
\centering
\small
\caption{معیارهای تفکیک‌شده برچسب‌ها در مجموعه‌داده \lr{DR}}
\label{tab:per-label-dr}
\begin{tabular}{llrrrr}
\hline
روش & برچسب & پشتیبانی & صحت & بازخوانی & \lr{F1-score} \\
\hline
\lr{Vanilla RAG} & خیر & ۱۸ & \lr{44.74\%} & \lr{94.44\%} & \lr{60.71\%} \\
\lr{Vanilla RAG} & بله & ۵۲ & \lr{96.88\%} & \lr{59.62\%} & \lr{73.81\%} \\
\lr{HippoRAG 2} & خیر & ۱۸ & \lr{45.95\%} & \lr{94.44\%} & \lr{61.82\%} \\
\lr{HippoRAG 2} & بله & ۵۲ & \lr{96.97\%} & \lr{61.54\%} & \lr{75.29\%} \\
\lr{CoT} & خیر & ۱۸ & \lr{56.67\%} & \lr{94.44\%} & \lr{70.83\%} \\
\lr{CoT} & بله & ۵۲ & \lr{97.50\%} & \lr{75.00\%} & \lr{84.78\%} \\
\lr{ToG-1} & خیر & ۱۸ & \lr{51.52\%} & \lr{94.44\%} & \lr{66.67\%} \\
\lr{ToG-1} & بله & ۵۲ & \lr{97.30\%} & \lr{69.23\%} & \lr{80.90\%} \\
\lr{ToG-2} & خیر & ۱۸ & \lr{51.52\%} & \lr{94.44\%} & \lr{66.67\%} \\
\lr{ToG-2} & بله & ۵۲ & \lr{97.30\%} & \lr{69.23\%} & \lr{80.90\%} \\
\hline
\end{tabular}
\end{table}

\begin{table}[ht]
\centering
\scriptsize
\caption{معیارهای تفکیک‌شده برچسب‌ها در مجموعه‌داده \lr{SWMH}}
\label{tab:per-label-swmh}
\begin{tabular}{llrrrr}
\hline
روش & برچسب & پشتیبانی & صحت & بازخوانی & \lr{F1-score} \\
\hline
\lr{Vanilla RAG} & اضطراب & ۱۴ & \lr{90.91\%} & \lr{71.43\%} & \lr{80.00\%} \\
\lr{Vanilla RAG} & اختلال دوقطبی & ۱۲ & \lr{100.00\%} & \lr{66.67\%} & \lr{80.00\%} \\
\lr{Vanilla RAG} & افسردگی & ۱۲ & \lr{47.37\%} & \lr{75.00\%} & \lr{58.06\%} \\
\lr{Vanilla RAG} & بدون اختلال روانی & ۱۴ & \lr{54.17\%} & \lr{92.86\%} & \lr{68.42\%} \\
\lr{Vanilla RAG} & خودکشی & ۱۳ & \lr{100.00\%} & \lr{23.08\%} & \lr{37.50\%} \\
\lr{HippoRAG 2} & اضطراب & ۱۴ & \lr{83.33\%} & \lr{71.43\%} & \lr{76.92\%} \\
\lr{HippoRAG 2} & اختلال دوقطبی & ۱۲ & \lr{100.00\%} & \lr{58.33\%} & \lr{73.68\%} \\
\lr{HippoRAG 2} & افسردگی & ۱۲ & \lr{44.44\%} & \lr{66.67\%} & \lr{53.33\%} \\
\lr{HippoRAG 2} & بدون اختلال روانی & ۱۴ & \lr{60.00\%} & \lr{85.71\%} & \lr{70.59\%} \\
\lr{HippoRAG 2} & خودکشی & ۱۳ & \lr{83.33\%} & \lr{38.46\%} & \lr{52.63\%} \\
\lr{CoT} & اضطراب & ۱۴ & \lr{70.59\%} & \lr{85.71\%} & \lr{77.42\%} \\
\lr{CoT} & اختلال دوقطبی & ۱۲ & \lr{100.00\%} & \lr{66.67\%} & \lr{80.00\%} \\
\lr{CoT} & افسردگی & ۱۲ & \lr{50.00\%} & \lr{75.00\%} & \lr{60.00\%} \\
\lr{CoT} & بدون اختلال روانی & ۱۴ & \lr{76.92\%} & \lr{71.43\%} & \lr{74.07\%} \\
\lr{CoT} & خودکشی & ۱۳ & \lr{77.78\%} & \lr{53.85\%} & \lr{63.64\%} \\
\lr{ToG-1} & اضطراب & ۱۴ & \lr{72.22\%} & \lr{92.86\%} & \lr{81.25\%} \\
\lr{ToG-1} & اختلال دوقطبی & ۱۲ & \lr{100.00\%} & \lr{83.33\%} & \lr{90.91\%} \\
\lr{ToG-1} & افسردگی & ۱۲ & \lr{43.75\%} & \lr{58.33\%} & \lr{50.00\%} \\
\lr{ToG-1} & بدون اختلال روانی & ۱۴ & \lr{88.89\%} & \lr{57.14\%} & \lr{69.57\%} \\
\lr{ToG-1} & خودکشی & ۱۳ & \lr{66.67\%} & \lr{61.54\%} & \lr{64.00\%} \\
\lr{ToG-2} & اضطراب & ۱۹ & \lr{100.00\%} & \lr{94.74\%} & \lr{97.30\%} \\
\lr{ToG-2} & اختلال دوقطبی & ۶ & \lr{100.00\%} & \lr{83.33\%} & \lr{90.91\%} \\
\lr{ToG-2} & افسردگی & ۲۰ & \lr{73.91\%} & \lr{85.00\%} & \lr{79.07\%} \\
\lr{ToG-2} & بدون اختلال روانی & ۱۱ & \lr{72.73\%} & \lr{72.73\%} & \lr{72.73\%} \\
\lr{ToG-2} & خودکشی & ۷ & \lr{66.67\%} & \lr{57.14\%} & \lr{61.54\%} \\
\hline
\end{tabular}
\end{table}

\newpage
\section{مقایسه توضیح‌پذیری روش‌ها}

برای ارزیابی توضیح‌های تولیدشده، دو دیدگاه مکمل به کار گرفته شد. نخست، شباهت معنایی توضیح تولیدشده با توضیح مرجع با استفاده از امتیازدهی مبتنی بر \lr{BERTScore} بررسی شد. دوم، ادعاهای بالینی موجود در توضیح‌ها با روش «ارزیابی سازگاری ادعاهای بالینی مبتنی بر گراف دانش» بررسی شدند؛ این روش از ایده‌های چارچوب \lr{GraphEval} الهام گرفته و میزان پشتیبانی ادعاهای استخراج‌شده از گراف دانش و منابع آن را اندازه‌گیری می‌کند \cite{sansford2024grapheval}. نتایج این دو ارزیابی در دو زیربخش بعدی ارائه می‌شوند.

\subsection{شباهت معنایی توضیح‌ها}

در این ارزیابی، توضیح تولیدشده و توضیح مرجع هر نمونه به بردار تبدیل و شباهت کسینوسی میان دو بردار محاسبه شد. این امتیاز یک تقریب مبتنی بر بردارِ کل توضیح است و با نسخه اصلی \lr{BERTScore} که در سطح نشانه‌ها محاسبه می‌شود تفاوت دارد. مقدار بیشتر نشان‌دهنده شباهت معنایی بالاتر با توضیح مرجع است.

\begin{table}[H]
\centering
\small
\caption{شباهت برداری توضیح‌های تولیدشده با توضیح مرجع}
\label{tab:explanation-embedding-similarity}
\begin{tabular}{lrrr}
\hline
روش & تعداد رکورد & میانگین امتیاز & میانه امتیاز \\
\hline
\lr{Vanilla RAG} & ۱۳۵ & \lr{0.7535} & \lr{0.7548} \\
\lr{HippoRAG 2} & ۱۳۵ & \lr{0.7451} & \lr{0.7551} \\
\lr{CoT} & ۱۳۵ & \lr{0.7774} & \lr{0.7832} \\
\lr{ToG-1} & ۱۳۵ & \lr{0.6611} & \lr{0.6751} \\
\lr{ToG-2} & ۱۳۵ & \lr{0.7180} & \lr{0.7232} \\
\hline
\end{tabular}
\end{table}

امتیاز \lr{BERTScore} به‌تنهایی معیار کاملی برای ارزیابی توضیح و استدلال مبتنی بر یک دانش یا توضیح مرجع نیست، زیرا عمدتاً شباهت سطحی دو متن را بر اساس واژه‌ها و موجودیت‌های مشابه اندازه‌گیری می‌کند. در نتیجه، ممکن است دو توضیح از نظر ادعاهای مطرح‌شده کاملاً متفاوت باشند، اما به دلیل شباهت واژگانی و معنایی امتیاز بالایی بگیرند. این محدودیت، ضرورت استفاده از ارزیابی‌های ریزدانه و ناظر بر ادعاهای توضیح را نشان می‌دهد.

\subsection{ارزیابی سازگاری ادعاهای بالینی مبتنی بر گراف دانش}

روش \lr{GraphEval} با هدف سنجش میزان وفاداری مدل زبانی به زمینه‌ای طراحی شده است که در ورودی در اختیار آن قرار می‌گیرد. هدف این پژوهش، ارزیابی یک مدل زبانی خاص از نظر وفاداری به ورودی آن نیست؛ بلکه می‌خواهیم توضیح‌های استدلالی روش‌های مختلف را بر اساس یک دانش مرجع، مانند کتاب \lr{DSM-5-TR}، ارزیابی کنیم. به همین دلیل، در روش «ارزیابی سازگاری ادعاهای بالینی مبتنی بر گراف دانش»، ابتدا توضیح مدل با کمک یک مدل زبانی به مجموعه‌ای از سه‌تایی‌های موجودیت، رابطه و موجودیت شکسته می‌شود. در همان دستور، ادعاها به دو گروه ادعاهای بالینی و ادعاهای مربوط به مورد یا نویسنده تقسیم می‌شوند؛ زیرا ممکن است مدل ادعا کند «پست‌کننده احساس نیاز به دیگران می‌کند» که ادعایی مربوط به مورد و نه یک ادعای پزشکی است، اما در جای دیگر بگوید «احساس نیاز به دیگران نشانه افسردگی است» که یک ادعای پزشکی محسوب می‌شود. ادعای دوم باید در برابر منبع مرجع، مانند کتاب \lr{DSM-5-TR}، بررسی شود تا مشخص شود آیا این منبع آن را پشتیبانی می‌کند یا نه.

پس از استخراج ادعاها، فقط ادعاهای بالینی در مجموعه واقعیت‌های \lr{Milvus} جست‌وجو می‌شوند و برای هر ادعا سه واقعیت مشابه بازگردانده می‌شود. سپس هر ادعا همراه با سه واقعیت بازیابی‌شده به مدل زبانی داده می‌شود تا میزان پشتیبانی واقعیت‌ها از ادعا را در بازه صفر تا ۱۰ امتیازدهی کند. در این گزارش، نرخ پشتیبانی برابر با درصد ادعاهایی است که امتیاز پشتیبانی آن‌ها بیشتر از ۵ بوده است. نتایج این امتیازدهی در جدول \ref{tab:clinical-claim-consistency} گزارش شده‌اند. این روش نیز مانند هر روش ارزیابی دیگری محدودیت‌هایی دارد که در بخش محدودیت‌های پژوهش بررسی خواهند شد.

\begin{table}[H]
\centering
\small
\caption{نتایج ارزیابی سازگاری ادعاهای بالینی با گراف دانش}
\label{tab:clinical-claim-consistency}
\begin{tabular}{lrrrrr}
\hline
منبع & تعداد نمونه & تعداد ادعا & ادعاهای پشتیبانی‌شده & نرخ پشتیبانی & میانگین امتیاز از ۱۰ \\
\hline
توضیح مرجع & ۱۳۵ & ۵۷۱ & ۳۸۳ & \lr{67.08\%} & \lr{6.21} \\
\lr{Vanilla RAG} & ۱۳۵ & ۶۶۷ & ۴۵۸ & \lr{68.67\%} & \lr{6.50} \\
\lr{HippoRAG 2} & ۱۳۵ & ۵۷۳ & ۴۰۴ & \lr{70.51\%} & \lr{6.79} \\
\lr{CoT} & ۱۳۵ & ۶۰۷ & ۴۲۴ & \lr{69.85\%} & \lr{6.63} \\
\lr{ToG-1} & ۱۳۵ & ۴۱۹ & ۳۳۹ & \lr{80.91\%} & \lr{7.87} \\
\lr{ToG-2} & ۱۳۵ & ۴۲۴ & ۳۰۶ & \lr{72.17\%} & \lr{6.78} \\
\hline
\end{tabular}
\end{table}

از نظر شباهت معنایی، \lr{CoT} بالاترین میانگین را دارد؛ اما از نظر سازگاری ادعاهای بالینی با گراف دانش، \lr{ToG-1} بیشترین نرخ پشتیبانی و بالاترین میانگین امتیاز را نشان می‌دهد. این نتیجه بیانگر آن است که شباهت به توضیح مرجع و اتکای قابل بررسی به شواهد گرافی دو جنبه متفاوت از توضیح‌پذیری هستند.

\section{مقایسه هزینه محاسباتی و زمان اجرا}

برای مقایسه هزینه محاسباتی، فقط میانگین مصرف منابع در هر نمونه گزارش شده است. معیارهای بررسی‌شده شامل میانگین نشانه‌های ورودی مدل زبانی، میانگین نشانه‌های خروجی، مجموع میانگین نشانه‌ها، میانگین تأخیر اجرای نمونه و میانگین تعداد فراخوانی‌های مدل زبانی است. در محاسبه میانگین توکن، خروجی‌هایی که مقدار ثبت‌شده آن‌ها صفر بوده است کنار گذاشته شدند؛ بنابراین، مقادیر جدول بر اساس اجراهای معتبر محاسبه شده‌اند.

\begin{table}[H]
\centering
\scriptsize
\caption{میانگین مصرف توکن، تأخیر و فراخوانی مدل زبانی در هر نمونه}
\label{tab:resource-usage}
\begin{tabular}{lrrrrr}
\hline
روش & میانگین توکن ورودی & میانگین توکن خروجی & میانگین کل توکن & میانگین تأخیر (میلی‌ثانیه) & میانگین فراخوانی مدل \\
\hline
\lr{Vanilla RAG} & \lr{2331.01} & \lr{157.59} & \lr{2488.59} & \lr{2032.70} & \lr{1.00} \\
\lr{HippoRAG 2} & \lr{5216.00} & \lr{543.55} & \lr{5759.55} & \lr{4568.80} & \lr{2.97} \\
\lr{CoT} & \lr{628.92} & \lr{137.55} & \lr{766.47} & \lr{1545.03} & \lr{1.00} \\
\lr{ToG-1} & \lr{35382.01} & \lr{10672.69} & \lr{46054.70} & \lr{430430.81} & \lr{45.90} \\
\lr{ToG-2} & \lr{9681.34} & \lr{603.96} & \lr{10285.30} & \lr{32362.88} & \lr{3.59} \\
\hline
\end{tabular}
\end{table}

بر اساس جدول \ref{tab:resource-usage}، روش \lr{CoT} کمترین مصرف توکن را دارد و \lr{Vanilla RAG} نیز با یک فراخوانی مدل در هر نمونه، هزینه و تأخیر محدودی ایجاد می‌کند. \lr{HippoRAG 2} به دلیل چند مرحله بازیابی و پالایش، هزینه بیشتری از \lr{Vanilla RAG} دارد. \lr{ToG-1} پرهزینه‌ترین روش است، زیرا در هر عمق چندین فراخوانی برای پالایش رابطه‌ها و امتیازدهی موجودیت‌ها انجام می‌دهد. \lr{ToG-2} با ادغام بخشی از تصمیم‌ها و استفاده از رتبه‌بندی متنی، نسبت به \lr{ToG-1} کاهش محسوسی در مصرف توکن و تأخیر نشان می‌دهد، هرچند همچنان از روش‌های مستقیم و بازیابی پایه پرهزینه‌تر است.

تأخیر گزارش‌شده برای \lr{ToG-1} و \lr{ToG-2} بر اساس اجرای نمونه محاسبه شده است و موازی‌سازی درون هر عمق را به‌صورت کامل فرض نمی‌کند. در این دو روش، فراخوانی‌های مربوط به چند رابطه یا چند موجودیت مستقل از یکدیگرند و می‌توان آن‌ها را هم‌زمان ارسال کرد؛ در نتیجه، زمان واقعی اجرای دسته‌ای می‌تواند از جمع زمان فراخوانی‌ها کمتر باشد. برای نمونه، اگر اجرای یک نمونه به‌صورت متوالی حدود ۲۰ دقیقه طول بکشد، با اجرای موازی فراخوانی‌های مستقل می‌توان زمان دیواری را تقریباً به ۵ دقیقه کاهش داد. این مقدار یک برآورد وابسته به ظرفیت هم‌زمانی سرویس، محدودیت نرخ درخواست و زمان کندترین فراخوانی است و به‌عنوان مقدار قطعی گزارش نمی‌شود.

\section{جمع‌بندی}

نتایج نشان دادند که عملکرد روش‌ها به مجموعه‌داده و نوع برچسب وابسته است و در این مقایسه، \lr{ToG-2} در نتایج کلی و \lr{CoT} در شباهت معنایی توضیح‌ها عملکرد برجسته‌ای داشتند. روش‌های گراف‌محور، به‌ویژه \lr{ToG-1}، در پشتیبانی ادعاهای بالینی با شواهد مرجع موفق‌تر بودند، اما هزینه محاسباتی و زمان اجرای بیشتری داشتند. در مقابل، روش‌های مستقیم و بازیابی پایه سریع‌تر و کم‌هزینه‌تر بودند و این فصل موازنه میان دقت، توضیح‌پذیری و هزینه اجرا را نشان داد.
